\section{Sharp consequences and the limits of uniformity}\label{sec:sharp}
\subsection{Three distinctions forced by retained information}
\begin{proposition}[Boundary scale and classwise physical time]\label{prop:boundarysharp}
The recurrence term in Theorem~\ref{thm:main} cannot in general be removed, even with deterministic bounded cycle length. The classwise ratios in \eqref{eq:classratio} cannot in general be replaced by a pooled ratio. Compactness of the unrestricted program space alone does not imply continuity of its average losses.
\end{proposition}
\begin{proof}
For $0<\eta<1/2$, take the retained boundary matrix
$P_\eta=\left(\begin{smallmatrix}1-\eta&\eta\\\eta&1-\eta\end{smallmatrix}\right)$,
initial label zero, and reward indicator of label one, constant during a deterministic two-call cycle. Over $n$ complete cycles the expected average is
\begin{equation}\label{eq:slow}
 {1\over2}-{1-(1-2\eta)^n\over4\eta n},
 \qquad h(P_\eta)={1\over2\eta}.
\end{equation}
Thus the error is of order $\min\{1,h(P_\eta)/n\}$. At $\eta=0$ the initial label is absorbing and the average is zero, whereas it is $1/2$ for every $\eta>0$. This proves both the necessity of a recurrence scale in a uniform bound and the possible failure of unrestricted continuity. The completely periodic swap matrix has a finite group inverse and is included in Theorem~\ref{thm:regularity}; period is not the obstruction.

For the second assertion, from an initial transient state enter either of two absorbing labels with probability $1/2$. Let their deterministic cycle lengths be two and four, and their per-call losses zero and one. The true average is $1/2$. The pooled ratio would be $(0+4)/(2+4)=2/3$. The difference is physical class weighting, not a matter of notation.

For a perturbation counterpart, change one of the two rates $\eta$ to $\eta+t$, $0<t\le\eta$. The stationary probability changes by $t/[2(2\eta+t)]$, whereas the one-cycle row total variation is $t$. Recurrence-sensitive amplification of order $h(P)t$ is therefore genuine. These finite matrices are witnesses for the general mechanism, not its proof.
\end{proof}

A physical reset also does not make learning across cycles vacuous. For an unknown fixed Bernoulli bias, a retained finite-state rule may accumulate evidence across independent acquisitions. Classical finite-memory hypothesis testing already treats this phenomenon \cite{HellmanCover}. In particular a sharpness example using a \emph{fixed} hidden bit that is revealed on ordinary short cycles cannot be transferred unchanged from a compulsory-reset protocol: a retained observer can learn that bit and use it on later erased cycles. The next construction uses a fresh scored coordinate and explicitly specifies the common task wire.

\subsection{A matched memory--calibration--defect slice}
A \emph{task-preserving} simulator below must pass through the environment-owned audit coordinate unchanged. It may transform visible reports and maintain private labels, but may not synthesize a replacement ground-truth audit. The target's load report must be completed before its first scored audit. This is the common-task causal interface being compared, not unrestricted Blackwell simulation with freely replaceable audit output.

\begin{theorem}[An attained nonreset joint resource law]\label{thm:joint}
Let a fresh vector $V$ at every physical preparation have a compact law $\pi$ with $q_M(\pi)\asymp M^{-2/d}$. Independently draw a fresh hidden fair bit $B$ and a lifetime $L$, with $L=1$ except on a marked rare event of probability $e\in(0,1/2]$, on which $L=\ell\ge1$. Preparation is charged and has zero score. At load the vector is visible. In the target, the bit is also visible; in the source it is replaced by a tagged erasure on rare cycles. Later visible holding reports contain neither a new vector nor information about the erased bit. At each of $L$ scored data calls the actual audit target is
$(V,u,B)$, where $u\in\{-\delta,\delta\}$ is one fixed inaccessible calibration parameter and $0\le\delta\le1$. Audits are not feedback. The observer is not overwritten at boundaries.

For task-preserving simulators the conditional completed-cycle deficiency of source to target is exactly
\begin{equation}\label{eq:attainedeps}
 \varepsilon=e/2.
\end{equation}
Put $\mu=1+\E L$ and $w=\E L/\mu$. Over \emph{all} stationary retained $M$-state observers and both fixed calibration alternatives, the unnormalized squared risk satisfies
\begin{align}\label{eq:jointbounds}
 wq_M(\pi)+w\delta^2+{e\ell\over4\mu}
 &\le \mathcal R_\av(M,\delta,e),\\
 \mathcal R_\av(M,\delta,e)
 &\le wq_J(\pi)+w\delta^2+{e\ell\over4\mu},
 \quad J=\lfloor M/3\rfloor,\quad M\ge3.\nonumber
\end{align}
The upper program uses at most $3J$ labels and has $h(P)\le2$; its preparation label is one of those existing labels because preparation is unscored and load/hold types are visible. The lower allows arbitrary retained controllers and fresh randomization, not merely the upper architecture.

For $0<p\le1$, choose $\ell=\lceil e^{-1/(1+p)}\rceil$. The family has a uniform $1+p$ cycle moment, and
\begin{equation}\label{eq:jointlaw}
 \mathcal R_\av(M,\delta,\varepsilon)
 \asymp M^{-2/d}+\delta^2+\varepsilon^{p/(1+p)}\qquad(M\ge3).
\end{equation}
All constants can be uniform in $e,\delta,M$ with the displayed preparation law and $p$ fixed. Dividing the task by the fixed bound $\operatorname{diam}(\supp\pi)^2+5$ gives a loss in $[0,1]$. The finite-acquisition comparison on a certified recurrence class is the \emph{signed} error of Theorem~\ref{thm:main}, not an additional positive learning penalty.
\end{theorem}
\begin{proof}
Conditional on any incoming observer or simulator label, the new bit is fair and independent of all past cycles and the vector. On a rare source load, no permitted information before the first audit reveals it. A guessed target bit agrees with the unaltered audit bit with probability at most $1/2$. In the target this agreement is certain. The measurable event ``rare and load bit equals first audit bit'' therefore gives total variation at least $e/2$ for the complete common-task law. Guessing a fresh fair bit on rare cycles and passing the true audit wire unchanged attains that distance; on ordinary cycles simply pass through the reported bit. Repeating the guess consistently on that cycle uses a finite simulator state, which is counted in a composed implementation. The audit deadline and identity wire are indispensable: if an independent artificial audit could be substituted, this would not be the same deficiency problem.

For any observer, its at most $M$ vector responses have pointwise squared error at least their nearest-center distortion. The vector is fresh and independent of lifetime and all incoming labels, so the acquired vector submeasure is $w\pi$. For the calibration lower, average the two fixed alternatives only as a minimax lower device. Their visible laws are identical, hence
$\tfrac12[(a-\delta)^2+(a+\delta)^2]=a^2+\delta^2\ge\delta^2$
at every scored call. The parameter is not redrawn during execution. On each rare cycle the fresh bit remains independent of the entire visible history throughout its holds, and any response has mean squared error at least $1/4$. Summing these three coordinate lower bounds gives the first line of \eqref{eq:jointbounds}. Retained memory cannot invalidate freshness of the current bit; that is exactly the difference from a fixed hidden-bit example.

For the upper, use $J$ vector centers crossed with bit outputs $0,1,1/2$, and calibration output zero. At load choose the nearest vector center and the reported bit, or $1/2$ after erasure; hold that label thereafter. The total count is $3J\le M$, and the already existing label at preparation incurs no score. Its next-cycle boundary label distribution is independent of its input label because every fresh load replaces it voluntarily. Thus $P=\Pi_P$ and $h(P)=\norm{I-P}_\infty\le2$. This proves the upper line.

Finally $\ell\le2e^{-1/(1+p)}$ and $\ell+1\le3e^{-1/(1+p)}$, so
$\E\tau^{1+p}\le2^{1+p}+3^{1+p}$.
Moreover $\mu$ is bounded above and below uniformly, $w\ge1/2$, and $e\ell\asymp e^{p/(1+p)}$. The comparison of $q_J$ and $q_M$ gives \eqref{eq:jointlaw}. Readouts may be restricted to the convex task ranges; the vector error is at most its squared diameter, calibration error at most $4\delta^2\le4$, and bit error at most one. This proves the normalization and its uniformity for all admissible programs.
\end{proof}

\begin{proposition}[Return-moment sharpness without memory erasure]\label{prop:tailsharp}
Uniformly over experiments with a bounded $1+p$ cycle moment, neither the $N^{-p}$ nor the $(1-\beta)^p$ order in Theorem~\ref{thm:main} can be improved. These are class-sharp orders, not lower bounds for every fixed experiment.
\end{proposition}
\begin{proof}
One retained label suffices. Let a cycle have one unscored preparation and either one zero-loss data call or $\ell$ unit-loss data calls; the long case has probability $e=\ell^{-(1+p)}$. The cycle moment is uniformly bounded and its long-run loss is $g=e\ell/[2+e(\ell-1)]\asymp\ell^{-p}$. Before $N$ calls there are at most $N$ possible cycle starts. The probability of any rare start is at most $Ne$, so $J_N\le Ne$. Choose $N=\lfloor c\ell\rfloor$ for a sufficiently small fixed $c>0$. Then $g-J_N\ge c'\ell^{-p}\asymp N^{-p}$. The lower uses an event probability, not a count of rewards per possible start.

For normalized discount, the probability of a rare start before call $t$ is at most $te$, whence
$J_\beta\le e/(1-\beta)$. Choose $(1-\beta)^{-1}=c\ell$ with the same sufficiently small $c$. This again gives $g-J_\beta\ge c'\ell^{-p}\asymp(1-\beta)^p$. There is no observer state to erase in this example, so it belongs to the retained class with $H=0$. The sharpness of recurrence-sensitive perturbation and of the tail exponent are separate, as the two examples show.
\end{proof}
